% \documentclass[aspectratio=169,notes]{beamer}
\documentclass[aspectratio=169]{beamer}
\usetheme[faculty=phil]{fibeamer}
\usepackage{polyglossia}
\setmainlanguage{english} %% main locale instead of `english`, you
%% can typeset the presentation in either Czech or Slovak,
%% respectively.
\setotherlanguages{russian} %% The additional keys allow
%%
%%   \begin{otherlanguage}{czech}   ... \end{otherlanguage}
%%   \begin{otherlanguage}{slovak}  ... \end{otherlanguage}
%%
%% These macros specify information about the presentation
\title[Artificial Intelligence Software (AIS)]{Artificial Intelligence Software, HW 4} %% that will be typeset on the
\subtitle{Data versioning by ClearML, Labeling tools  
\\ \  \\ \  } %% title page.
\author{Oleg Bulichev}
%% These additional packages are used within the document:
\usepackage{ragged2e}  % `\justifying` text
\usepackage{booktabs}  % Tables
\usepackage{tabularx}
\usepackage{tikz}      % Diagrams
\usetikzlibrary{calc, shapes, backgrounds}
\usetikzlibrary{decorations.pathreplacing,calligraphy,calc,graphs}
\usepackage{amsmath, amssymb}
\usepackage{url}       % `\url`s
\usepackage{listings}  % Code listings
% \usepackage{subfigure}
\usepackage{floatrow}
\usepackage{subcaption}
\usepackage{mathtools}
\usepackage{todonotes}
\usepackage{fontspec}
\usepackage{multicol}
\usepackage{pdfpages}
\usepackage{wrapfig}
\usepackage{qrcode}
\usepackage{animate}
\usepackage{booktabs}
\usepackage{multirow}
% \usepackage{graphicx}
\usepackage{colortbl}

\graphicspath{{resources/}}
\frenchspacing

\setbeamertemplate{caption}[numbered]
\usetikzlibrary{graphs}

% \usepackage[backend=biber,style=ieee,autocite=footnote]{biblatex}
% \addbibresource{biblio.bib}
% \DefineBibliographyStrings{english}{%
%   bibliography = {References},}

\newcommand{\oleg}[2][] {\todo[color=red, #1] {OLEG:\\ #2}}
\newcommand{\fbckg}[1]{\usebackgroundtemplate{\includegraphics[width=\paperwidth]{#1}}}%frame background

\usepackage[framemethod=TikZ]{mdframed}
\newcommand{\dbox}[1]{
\begin{mdframed}[roundcorner=3pt, backgroundcolor=yellow, linewidth=0]
\vspace{1mm}
{#1}
\vspace{1mm}
\end{mdframed}
}

\begin{document}
\setlength{\abovedisplayskip}{0pt}
\setlength{\belowdisplayskip}{0pt}
\setlength{\abovedisplayshortskip}{0pt}
\setlength{\belowdisplayshortskip}{0pt}

\fbckg{fibeamer/figs/title_page.png}
\frame[c]{\setcounter{framenumber}{0}
    \usebeamerfont{title}%
    \usebeamercolor[fg]{title}%
    \begin{minipage}[b][6.5\baselineskip][b]{\textwidth}%
        \textcolor{black}{\raggedright\inserttitle}
    \end{minipage}
    % \vskip-1.5\baselineskip

    \usebeamerfont{subtitle}%
    \usebeamercolor[fg]{framesubtitle}%
    \begin{minipage}[b][3\baselineskip][b]{\textwidth}
        \raggedright%
        \insertsubtitle%
    \end{minipage}
    \vskip.25\baselineskip
}
%   \frame[c]{\maketitle}

\fbckg{fibeamer/figs/common.png}

\note{\scriptsize \begin{itemize}
        \item \
    \end{itemize}}

\begin{frame}[t]{Goal of the Assignment}
\vspace*{-0.5cm}

\textbf{Short description:} The task is to learn the full pipeline of preparing
datasets for computer vision tasks.

\begin{itemize}
\item Choose either detection or segmentation problem.
\item Find and analyze a public dataset.
\item Collect own image data.
\item Annotate images.
\item Add collected and public datasets to ClearML.
\item Merge datasets.
\item Split dataset into train/val/test.
\item Export the dataset in a format compatible with YOLO models
\end{itemize}

The final result should be a structured dataset suitable for training
a neural network.

\end{frame}

\begin{frame}[t]{Technology Stack}

The following tools should be used during the assignment.

\textbf{Dataset management}

\begin{itemize}
\item ClearML Dataset
\end{itemize}

\textbf{Annotation tools}

\begin{itemize}
\item CVAT
\item Roboflow
\item Supervisely
\end{itemize}

\textbf{Pre-annotation}

\begin{itemize}
\item Segment Anything Model (SAM) or others
\end{itemize}

You may use either local installations or cloud versions of these tools.

\end{frame}

\begin{frame}[t]{Task Selection}

Tasks are listed in order of increasing complexity:

\begin{enumerate}
\item Object Detection
\item Image Segmentation
\end{enumerate}

Annotation type depends on the selected task:

\begin{itemize}
\item Detection — bounding boxes
\item Segmentation — masks or polygons
\end{itemize}

\end{frame}

\begin{frame}[t]{Step 1: Public Dataset}
\vspace*{-0.5cm}
Find a publicly available dataset related to your selected task.

\begin{itemize}
\item Kaggle
\item Hugging Face Datasets
\item PapersWithCode
\item Roboflow Universe
\item Supervisely Datasets
\end{itemize}

In your report include:

\begin{itemize}
\item dataset name with the source link
\item short description
\item number of images
\end{itemize}

\end{frame}

\begin{frame}[t]{Step 2: Collect Your Own Images}
\vspace*{-0.5cm}
You must collect your own dataset. The goal is to build a realistic and diverse dataset.

Requirements:

\begin{itemize}
\item at least 100 images
\item images must be taken by you
\item dataset must contain variation
\end{itemize}

Ensure diversity in:

\begin{itemize}
\item camera angles
\item lighting conditions
\item backgrounds
\item scene composition
\end{itemize}



\end{frame}

\begin{frame}[t]{Step 3: Image Annotation}
\vspace*{-0.5cm}
All collected images must be annotated.

Choose an annotation tool:

\begin{itemize}
\item CVAT
\item Roboflow
\item Supervisely
\end{itemize}

Annotation type depends on the selected task:

\begin{itemize}
\item Detection — bounding boxes
\item Segmentation — masks or polygons
\end{itemize}

Ensure consistent annotation across all images.

\end{frame}


\begin{frame}[t, fragile]{Step 4: Create ClearML Dataset}
\vspace*{-0.5cm}

Create a ClearML Dataset for both the public and collected datasets. And fuse them together.

To sum up you should have three datasets in ClearML.
\end{frame}

\begin{frame}[t]{Step 5: Dataset Split}

Split the dataset into three parts:

\begin{itemize}
\item training set
\item validation set
\item test set
\end{itemize}

Recommended split:

\begin{itemize}
\item 70\% training
\item 20\% validation
\item 10\% testing
\end{itemize}

Ensure that all subsets contain representative samples.
\end{frame}




\fbckg{fibeamer/figs/last_page.png}
\frame[plain]{}

\end{document}